\documentclass[11pt]{article}
\usepackage[letterpaper, left=0.7in, right=0.7in, top=0.75in, bottom=0.75in,
  headheight=14pt, headsep=0.2in, footskip=0.4in]{geometry}
\usepackage[T1]{fontenc}
\usepackage[default]{sourcesanspro}
\usepackage{inconsolata}
\usepackage{microtype}
\usepackage{fancyhdr}
\usepackage{array}
\usepackage{tabularx}
\usepackage{booktabs}
\usepackage{enumitem}
\usepackage{titlesec}
\usepackage{amsmath}
\usepackage{amssymb}
\usepackage{xcolor}
\usepackage{needspace}

\pagestyle{fancy}
\fancyhf{}
\fancyhead[L]{Week 4 / Stars and wheels / Adult guide}
\fancyfoot[L]{Bellingham Math Circle / Week 4 / F04-FAC-v4}
\fancyfoot[R]{\thepage}
\renewcommand{\headrulewidth}{0pt}
\setlength{\parindent}{0pt}
\setlength{\parskip}{3pt}
\setlist{nosep, leftmargin=1.3em, topsep=1pt}

\titleformat{\section}{\Large\bfseries}{\thesection.}{0.45em}{}
\titlespacing*{\section}{0pt}{10pt}{3pt}
\titleformat{\subsection}{\large\bfseries}{}{0pt}{}
\titlespacing*{\subsection}{0pt}{7pt}{1pt}

\newcommand{\lbl}[1]{\colorbox{black!10}{\footnotesize\strut #1}}
\newcommand{\core}{\lbl{core}}
\newcommand{\further}{\lbl{further}}
\newcommand{\ext}{\lbl{extension}}
% \pb{number}{tag}{what the children do}
\newcommand{\pb}[3]{\par\needspace{4\baselineskip}\vspace{3pt}\noindent\textbf{Problem #1}\enspace #2\enspace #3\par}
\newcommand{\ans}{\textbf{Answer.}\ }
\newcommand{\say}{\textbf{Say.}\ }
\newenvironment{hints}{\begin{enumerate}[label={\textit{Hint \arabic*.}}, leftmargin=*, align=left]}{\end{enumerate}}
\newcolumntype{L}{>{\raggedright\arraybackslash}X}
\renewcommand{\arraystretch}{1.15}

\begin{document}

% =====================================================================================
\section{The idea of the week}

Put dots evenly around a circle. A \emph{hop of 3} means: move 3 dots clockwise and draw a straight line from
where you were to where you land. Keep making the same hop until you are back where you began. Sometimes the line
reaches every dot and draws a star. Sometimes it comes back early; then you start again at a dot with no line,
with the same hop, until every dot has a line.

The question all week: when does one hop reach every dot, and how many starts does it take? The answer: the
number of starts is the biggest number that goes evenly into both the number of dots and the hop. On 12 dots,
hop 8 needs 4 starts (four triangles), because 4 goes into 12 and into 8 and no bigger number does. On 7 dots
every hop reaches every dot.

The code wheel is the same idea with 26 letters. Turning the inner wheel 3 places moves every letter 3 along,
which is a hop of 3 on a ring of 26. Two turns in a row make one turn, turning back undoes a turn, and since
there are only 25 real turns, a code can be cracked by trying them. The K--1 table does the same with a ring of
six pictures and no letters.

\textbf{What a good session looks like}
\begin{itemize}
\item \textbf{K--1.} Every child hops a counter correctly (the first count lands on the next dot, not the one
  the counter is on) and draws hops on the small rings. Each child can say, for at least one ring, which hops
  landed on every dot and which came back early. A very good session reaches the ring of 7 or the pictures.
\item \textbf{2--3.} Every child draws hops on 8, 10 or 12 dots and counts the starts, guesses before drawing,
  and also codes and decodes a message on their own wheel with a partner. The best sign is a guess with a reason:
  ``hop 3 on 12 only lands on 0, 3, 6 and 9''.
\item \textbf{4--5.} The children state a rule for the number of pieces, test it on the big numbers of
  Problem 3, and try to say why it works. They crack codes from clues and find that two settings make one.
\end{itemize}

% =====================================================================================
\section{Materials and preparation}

Eleven children: K--1 four (two pairs), 2--3 four (two pairs), 4--5 three (a pair and a referee).

{\small\begin{tabularx}{\textwidth}{@{}l c c c c L@{}}
\toprule
Item & K--1 & 2--3 & 4--5 & Total & What for \\
\midrule
Counters & 20 & 20 & 20 & 61 & K--1: per pair, 1 to hop and 7 of another colour to mark landings, plus spares.
  2--3 and 4--5: equal-pile explanations. One more for the launch. \\
Coloured pencils, 6+ colours & 2 sets & 2 sets & 2 sets & 6 sets & One set per pair (4--5: one for the pair, one for the referee); a new colour for each start. \\
Pencils with erasers & 5 & 5 & 4 & 14 & One per child and adult. \\
Rulers, 6-inch & 2 & 4 & 3 & 9 & Straight lines on the small rings. K--1 rulers are for first graders who want one. \\
Scissors & -- & 4 & 3 & 8 & Children cut their own wheels; one pair for the adults' spare wheels. \\
Brass fasteners & -- & 4 & 3 & 10 & One per wheel: 7 for the children, 3 for spare wheels. Bring all 20; legs bend. \\
Small whiteboards, markers & -- & 1 & 1 & 3 & Trying settings; one more, pre-drawn, for the launch. \\
Scrap paper & -- & 10 & 10 & 20 & Messages and fallback code games. \\
\bottomrule
\end{tabularx}}

\medskip
\textbf{What to print.} Single-sided, US Letter, at 100\% (``Actual size'', not ``Fit to page''). Every child
gets their own copy, because every child draws.

{\small\begin{tabularx}{\textwidth}{@{}l l c L@{}}
\toprule
Table & Pages & Copies & How to hand out \\
\midrule
K--1 & \texttt{k-1.pdf} pp.\ 1--10 & 6 each & 4 children, 1 adult, 1 spare. One page at a time; keep the rest in order with you. \\
2--3 & \texttt{grades-2-3.pdf} pp.\ 1--7 & 6 each & 4 children, 1 adult, 1 spare. Staple pp.\ 1--7. \\
2--3 & \texttt{grades-2-3.pdf} pp.\ 8--9 (wheels) & 6 each & 4 children, 2 for spare wheels. Card stock if the printer takes it. \\
4--5 & \texttt{grades-4-5.pdf} pp.\ 1--9 & 5 each & 3 children, 1 adult, 1 spare. Staple pp.\ 1--9. \\
4--5 & \texttt{grades-4-5.pdf} pp.\ 10--11 (wheels) & 4 each & 3 children, 1 for a spare wheel. \\
\bottomrule
\end{tabularx}}

\medskip
\textbf{Before the session}
\begin{enumerate}
\item \textbf{Check the scale.} The outer wheel must measure 5.5 inches across and the inner wheel 4 inches. If
  not, the printer shrank the page: reprint at actual size.
\item \textbf{Check the counters on the K--1 dots.} The dots on the big rings (K--1 pp.\ 1--3, 6, 7) are
  0.8 inch (20 mm) across. Put a real counter on a printed page 1: it should sit on one dot without touching the
  next. A 3/4-inch counter fits inside a dot; a 1-inch counter overhangs a little and is fine if you can still see
  which dot it is on. If the counters are too big, use pennies. On the picture rings (pp.\ 8--10) a counter
  covers the picture, so children set it just outside the picture's circle.
\item \textbf{Pierce the wheel centres.} Stack the outer-wheel pages and push a pushpin through the centre cross;
  do the same with the inner-wheel pages. Children then only need to push the fastener through.
\item \textbf{Make 3 spare wheels} (2 for the 2--3 table, 1 for 4--5). Cut both circles on the line, put the
  inner wheel on top, push a fastener through both centres, and open the legs flat on the back, not tight.
  \textbf{Test one:} turn the inner A to the outer D; the inner C, A, T must sit beside the outer F, D, W, the
  inner wheel must turn freely, and every outer letter must show around the inner wheel.
\item \textbf{Cutting time.} Children cut their own wheels when they reach Problem 4. Cutting both circles and
  fastening takes a third grader about 10 minutes and a fourth or fifth grader about 5--7 minutes. Rough cutting
  is fine as long as no letter is cut; cut the inner wheel on its line or just inside it. A child still cutting
  after 10 minutes gets a spare wheel and finishes their own at home.
\item \textbf{Launch board.} On a whiteboard draw a ring of 5 dots as large as the board allows, each dot about
  1 inch across. Fill in the top dot and draw a clockwise arrow inside the ring. Have 1 counter.
\item \textbf{K--1 kits.} For each pair, a cup with 1 counter of one colour (the hopper) and 7 of another
  colour (markers). K--1 gets no letter wheel: its secret hops use the printed ring of pictures (pp.\ 8--10).
\end{enumerate}

% =====================================================================================
\section{The hour}

{\small\begin{tabularx}{\textwidth}{@{}l l L@{}}
\toprule
0:00--0:05 & Running around & \\
0:05--0:08 & Launch & Everyone stands around one table; the organizer leads. \\
0:08--0:48 & Tables & About 40 minutes. 2--3 and 4--5 should reach the wheel (Problem 4) by about 0:28. \\
0:48--0:55 & Sharing & One question, one answer from each table. \\
\bottomrule
\end{tabularx}}

\subsection{Launch: one hop, 2--3 minutes}
The organizer holds the launch board (ring of 5, black dot at the top) and one counter.
\begin{enumerate}
\item Adult puts the counter on the black dot. Says: \emph{``Five dots. A hop of 2 moves the counter two dots the
  way the arrow points.''}
\item Adult hops once, touching each dot: \emph{``One, two.''} Draws a straight line from the black dot to the
  dot where the counter landed.
\item Adult makes one wrong move on purpose: counts the dot the counter is sitting on as ``one'' and moves only one
  dot. Asks: \emph{``Was that a hop of 2?''} The children say no; the adult moves it back.
\item Says: \emph{``Who will do the next hop?''} One child moves the counter, touching each dot and saying ``one,
  two''. A legal move goes exactly two dots clockwise, with ``one'' on the next dot. The adult draws the line. The
  child keeps going until the counter is back on the black dot. The drawing is a five-pointed star.
\item Says: \emph{``Every dot got a line. Does that always happen? At your tables you will find out with other
  rings and other hops. The older tables will also make hops on a wheel of letters, to make secret codes.''}
\end{enumerate}

\subsection{Pairs at each table}
\begin{itemize}
\item \textbf{K--1, two pairs:} a kindergartner with a first grader. One child hops the counter and counts aloud;
  the partner puts a marker on each dot it lands on and says ``back'' when it reaches the black dot. Both draw the
  hop on their own small ring. Swap jobs for every new hop.
\item \textbf{2--3, two pairs.} On ring pages one child draws while the partner counts each hop aloud on the dots
  and checks that the line ends on the right dot; swap every ring. With wheels each child codes on their own wheel;
  in Problem 5 partners code for each other.
\item \textbf{4--5, a pair and a referee,} rotating at each new problem. On ring pages the referee chooses which
  rings each player draws and checks the lines and the count. With codes the referee writes a code with a secret
  setting and the pair cracks it.
\end{itemize}

\subsection{When attention runs out}
\begin{itemize}
\item \textbf{K--1: pass the counter.} Stand around the table: 4 children and the adult make 5 people. Pass a
  counter to the person 2 places along, everyone saying ``one, two''. Does everyone get it before it comes back?
  (Yes.) Then the adult steps out, leaving 4, and they pass by 2 again: two children never get it. Ask which pass
  would reach everyone (1 or 3).
\item \textbf{2--3: secret names.} One child writes an animal's name in code with a secret setting; the partner
  cracks it with their wheel; swap. Before wheels are made: colour each start of a finished drawing a different
  colour and count the colours.
\item \textbf{4--5: crack race.} The referee writes one word in code with a secret setting; the other two race to
  crack it; the winner is the next referee.
\end{itemize}

\subsection{Sharing question}
\emph{``Name a ring and a hop where you had to start again. Could you have told before you drew it?''}

% =====================================================================================
\section{Table by table}

\subsection{K--1 (parent volunteer)}
Read each problem aloud once, then show it with the counter. The children do not read; the rings and boxes carry
the task. ``Lands on every dot'' means every dot gets a marker before the counter is back on the black dot. In the
drawings, a line goes from the dot the counter left to the dot where it landed. Counting starts on the next dot:
on 4 dots, hop 3 from the top lands on the dot to the left of the top.

\textbf{Core:} Problems 1--4. \textbf{Further:} 8, then 6--7, 5, 9, 10. If drawing gets tiring after Problem 3
or 4, switch to Problem 8 (pictures). \textbf{If a child is struggling:} the child hops and you draw the lines;
do hops 1 and 2 on each ring; skip Problems 5, 9 and 10.

\pb{1--3}{\core}{Rings of 4, 5 and 6 dots, hops 1, 2, 3. Hop the counter on the big ring, draw each hop on its
small ring, circle the hops that land on every dot.}
\ans Circle hops 1 and 3 on 4 dots; all three on 5 dots; only hop 1 on 6 dots.

{\small\begin{tabularx}{\textwidth}{@{}l L L L@{}}
\toprule
& hop 1 & hop 2 & hop 3 \\
\midrule
4 dots & all 4 dots: a square \checkmark & 2 dots: one line, there and back & all 4: same square, other way \checkmark \\
5 dots & all 5: a pentagon \checkmark & all 5: a five-pointed star \checkmark & all 5: the same star the other way \checkmark \\
6 dots & all 6: a hexagon \checkmark & 3 dots: a triangle & 2 dots: one line across \\
\bottomrule
\end{tabularx}}
\begin{hints}
\item ``Put your finger on the next dot and say one. Where does the last number land?''
\item ``Put a marker on every dot the counter lands on. When it is back on the black dot, is any dot bare?''
\item ``Count the markers. How many dots are on the ring?''
\end{hints}
\say ``On 6 dots, hop 2 came back after 3 hops. Why so soon?'' (2, 4, 6: three hops make one trip round.)

\pb{4}{\core}{Draw the hops; when you are back and a dot has no line, start again there with the same hop.
Write how many times you started.}
\ans 6 dots, hop 2: \textbf{2} (two triangles make a six-pointed star). 6 dots, hop 3: \textbf{3} (three lines
through the middle). 8 dots, hop 2: \textbf{2} (two squares). 8 dots, hop 3: \textbf{1} (an eight-pointed star).
The first start counts: a child who writes 1, 2, 1, 0 is counting only the restarts. Ask ``How many times did you
put your pencil down to start?''
\begin{hints}
\item ``You are back at the black dot. Point to a dot with no line.''
\item ``Start there with the same hop, in a new colour.''
\item ``Count the colours.''
\end{hints}

\pb{5}{\further}{Find every hop from 1 to 5 that draws each picture. The four small rings are for trying.}
\ans Pentagon: hops \textbf{1, 4}. Five-pointed star: \textbf{2, 3}. Hexagon: \textbf{1, 5}. Six-pointed star:
\textbf{2, 4}. Three crossing lines: \textbf{3}. (Hop 5 on 5 dots does not move.)
\begin{hints}
\item ``Count the dots of this picture. Which ring has that many?''
\item ``Look back at your drawings on pages 2--4. Which one matches?''
\item ``Try hop 4 on a small ring of 5 dots. Which picture does it make?'' (Hop 4 on 5 dots is one step backwards.)
\end{hints}

\pb{6--7}{\further}{Ring of 7 dots: hops 1, 2, 3, then hops 4, 5, 6.}
\ans Every hop lands on every dot: circle all six. Hops 1 and 6 draw the seven-sided shape, hops 2 and 5 one
seven-pointed star, hops 3 and 4 a pointier seven-pointed star.
\begin{hints}
\item Markers again: ``Is any dot bare when it is back?''
\item ``On 6 dots some hops missed dots. Does any hop miss on 7?''
\item ``Put hop 6 next to hop 1. Same picture?''
\end{hints}
\say ``Seven counters cannot be put into equal piles, unless each pile is just one. That is why no hop comes back
early.''

\pb{8}{\further}{Hop 2 on the ring of pictures, row by row, until the top row comes back.}
\ans Clockwise from the top, the ring is sun, moon, heart, tree, fish, house. The rows are
\textbf{heart, house, sun}; then \textbf{fish, moon, heart}; then \textbf{sun, tree, fish} (the top row again).
The last two rows of boxes stay empty.
\begin{hints}
\item ``Put a counter beside the sun on the ring and hop it 2. Which picture is it next to?''
\item ``Draw that picture in the box under the sun. Now do the tree.''
\item ``Is the new row the same as the top row? If not, keep going.''
\end{hints}

\pb{9}{\further}{Which hop changes the picture back?}
\ans Hop 1 is undone by hop \textbf{5}, hop 2 by \textbf{4}, hop 3 by \textbf{3}, hop 4 by \textbf{2}. Each
pair adds to 6: once round the ring.
\begin{hints}
\item ``Put the counter on the middle picture. Hop one place at a time until you reach the picture on the right.
  Count the hops.''
\item ``Write that number in the box.''
\item ``Look at hop 1 and its answer, and hop 2 and its answer. What do they add up to?''
\end{hints}

\pb{10}{\further}{A secret hop changed the sun into the tree. Draw what it does to four pictures.}
\ans The secret hop is 3. Moon $\to$ \textbf{fish}, heart $\to$ \textbf{house}, fish $\to$ \textbf{moon},
house $\to$ \textbf{heart}.
\begin{hints}
\item ``Put the counter on the sun. Hop one place at a time to the tree and count.''
\item ``Start on the moon and make the same number of hops.''
\end{hints}

% -------------------------------------------------------------------------------------
\subsection{Grades 2--3 (mathematician)}
The rule is printed at the top of page 1. \textbf{Core:} Problems 1, 3, 4, 5, and 2 as time allows.
\textbf{Further:} 6, 7, 8, 9. Move every child to Problem 4 (the wheel) by about 20 minutes into table time, even
with Problems 2--3 unfinished; they can return to them. \textbf{If a child is struggling:} on Problem 2 draw only
hops 2 and 5; skip the guesses on Problem 3; decode two of the codes in Problem 4.

\pb{1, 2, 3, 6}{\core\ (6: \further)}{Draw each hop and count the starts; on Problems 3 and 6, guess first.}

{\small\begin{tabularx}{\textwidth}{@{}l l c L@{}}
\toprule
Problem & Ring, hop & Starts & What appears \\
\midrule
1 & 8 dots: hop 1, 2, 3, 4 & 1, 2, 1, 4 & octagon; two squares; eight-pointed star; four lines through the middle \\
2 & 10 dots: hop 2, 3, 4, 5 & 2, 1, 2, 5 & two pentagons; ten-pointed star; two five-pointed stars; five lines through the middle \\
3 & 12 dots: hop 2, 3, 4, 5, 6 & 2, 3, 4, 1, 6 & two hexagons; three squares; four triangles; twelve-pointed star; six lines through the middle \\
6 & 9 dots, hop 2 & 1 & nine-pointed star \\
  & 15 dots, hop 5 & 5 & five triangles \\
  & 16 dots, hop 6 & 2 & two eight-pointed stars \\
  & 12 dots, hop 8 & 4 & four triangles, the same picture as hop 4 drawn the other way \\
\bottomrule
\end{tabularx}}
\begin{hints}
\item ``Put your pencil on a dot and count the hop aloud, touching each dot.''
\item ``You are back where you started. Is there a dot with no line? Start there in a new colour.''
\item ``How many dots did one start reach? How many groups that size make all the dots?'' (Problem 6, 12 dots
  hop 8: ``Count the other way round. Which smaller hop lands on the same dot?'')
\end{hints}
On Problem 3 many children guess ``starts = hop''. It is right for hops 2, 3, 4, 6 (they go evenly into 12) and
wrong for hop 5. Ask why hop 5 is different.

\pb{4}{\core}{Cut out and join the wheel; with the wheel set to D, find what each code came from.}
\ans VWDU = \textbf{STAR}; IRA = \textbf{FOX}; ZKHHO = \textbf{WHEEL}; SHQFLO = \textbf{PENCIL};
L FDQ GUDZ D VWDU = \textbf{I CAN DRAW A STAR}. To decode, find the code letter on the outer wheel and read the
inner letter beside it. The usual mistake is to go inner to outer, which codes it again (VWDU becomes YZGX).
\begin{hints}
\item ``Is your inner A beside the outer D? Check that CAT gives FDW.''
\item ``Find V on the outer wheel. Which inner letter is beside it?''
\item ``You made code by going from inner to outer. Which way undoes it?''
\end{hints}

\pb{5}{\core}{Write a message of five or more words in code for your partner; swap and decode.}
\ans No fixed answer. Check that the partner's decoding gives back the message. Nonsense usually means a wrong
setting or a wheel that slipped.
\begin{hints}
\item ``Find each letter on the inner wheel and write the outer letter beside it.''
\item ``Hold the wheel still with your thumb while you code.''
\item ``Decode the first word together and compare.''
\end{hints}

\pb{7}{\further}{On rings of 9, 16, 18, 15 dots find a hop bigger than 3 with exactly 3 starts, or explain why
there is none.}
\ans All hops smaller than the ring: 9 dots \textbf{6}; 16 dots \textbf{none}; 18 dots \textbf{15} only (6, 9, 12
give 6, 9, 6 starts); 15 dots \textbf{6, 9 or 12}.
\textbf{Why not 16, with counters:} every start reaches the same number of dots, and the starts together use all
16 dots. Three starts would put 16 counters into 3 equal piles, but 16 counters make 5, 5, 5 with 1 left over.
\begin{hints}
\item ``Which hop gave 3 starts on 12 dots, in Problem 3?''
\item ``On 9 dots, hop 3 gives 3 starts. Which bigger hop lands on the same dots, going the other way?''
\item ``For 16 dots: put 16 counters into 3 equal piles.''
\end{hints}

\pb{8}{\further}{Crack two codes made with unknown settings.}
\ans Setting H: \textbf{A STAR CAN HAVE TEN POINTS}. Setting Q: \textbf{MEET ME BY THE BIG TREE}. The one-letter
word H must be A or I; A $\to$ H is setting H. In the second message CUUJ and JHUU have double letters; setting
Q codes E as U, so UU decodes to EE.
\begin{hints}
\item ``Which words have one letter? If H came from A, where must the wheel be set?''
\item ``Set it there and try the next word. Is it a word?''
\item ``CUUJ has a double letter in the middle. Think of four-letter words like that, and try the setting.''
\end{hints}

\pb{9}{\further}{Second setting that gives the word back.}
\ans D $\to$ \textbf{X}; H $\to$ \textbf{T}; N $\to$ \textbf{N}; W $\to$ \textbf{E}. The two turns add up to one
full turn of 26 (D is 3 places, X is 23).
\begin{hints}
\item ``Set the wheel to D and code A. Which setting takes that letter back to A?''
\item ``How many places did D move A? How many more to get all the way round?''
\item ``Why is N its own undo?'' (13 and 13 make 26.)
\end{hints}

% -------------------------------------------------------------------------------------
\subsection{Grades 4--5 (organizer)}
\textbf{Core:} Problems 1, 2, 3, 4, 6, 7. \textbf{Further:} 5, 8, 9, 10, 11; 12 is the extension. ``Pieces''
means starts: the page-1 rule says so, and the pieces of a drawing cross each other, so counting connected parts
would give 1 every time. \textbf{If a child is struggling:} draw only the first row of Problem 2 and go on to
Problem 4.

\pb{1--3}{\core}{Count the pieces: 12 dots with hops 1--6; six mixed rings; then seven cases without drawing.}

{\small\begin{tabularx}{\textwidth}{@{}l l c L@{}}
\toprule
Problem & Ring, hop & Pieces & What appears \\
\midrule
1 & 12 dots: hop 1, 2, 3, 4, 5, 6 & 1, 2, 3, 4, 1, 6 & 12-gon; 2 hexagons; 3 squares; 4 triangles; 12-pointed star \{12/5\}; 6 lines through the middle \\
2 & 10 dots, hop 4 & 2 & two five-pointed stars \\
  & 9 dots, hop 6 & 3 & three triangles \\
  & 15 dots, hop 10 & 5 & five triangles \\
  & 16 dots, hop 12 & 4 & four squares \\
  & 18 dots, hop 8 & 2 & two nine-pointed stars \{9/4\} \\
  & 24 dots, hop 9 & 3 & three eight-pointed stars \{8/3\} \\
3 & 20/8, 24/10, 30/12, 36/27 & 4, 2, 6, 9 & 20/8: four five-pointed stars; 24/10: two \{12/5\} stars \\
  & 17/5, 100/35, 60/45 & 1, 5, 15 & \\
\bottomrule
\end{tabularx}}
\begin{hints}
\item ``Colour each piece differently and count the colours.''
\item ``Number the dots 0, 1, 2, \dots{} Which numbers does the first piece land on?''
\item ``What does the number of pieces have to do with the dots and the hop? Try the biggest number that goes
  into both.''
\end{hints}

\pb{4}{\core}{Make the wheel; crack four codes; find two words for MROOB.}
\ans Setting J: \textbf{I DREW A STAR WITH TWELVE POINTS}. Setting S: \textbf{WE MEET AT NOON BY THE BIG TREE}.
Setting X: \textbf{BALLOON}. MROOB: \textbf{CHEER} (setting K) and \textbf{JOLLY} (setting D); the only other
dictionary word is DIFFS (setting J). Clues: a one-letter word is A or I; YXIILLK has two double letters.
\begin{hints}
\item ``Which words have one letter? If R came from I, where is the wheel set?''
\item ``Look at the double letters. Which seven-letter word has two pairs of doubles?''
\item ``For MROOB, try every setting and write all 25 results.''
\end{hints}

\pb{5}{\further}{Code a message of six or more words with a secret setting; a partner cracks it.}
\ans No fixed answer; the referee checks. Good as a fallback.
\begin{hints}
\item ``Find each letter on the inner wheel and write the outer letter beside it.''
\item Referee: ``Decode the first word yourself and compare.''
\end{hints}

\pb{6}{\core}{Two settings in a row: which single setting does the same?}
\ans H, K $\to$ \textbf{R}; D, F $\to$ \textbf{I}; T, M $\to$ \textbf{F}; P, L $\to$ \textbf{A}; missing second
setting after H for A: \textbf{T}. Add the places (A = 0, B = 1, \dots), going round past Z: T then M is
19 + 12 = 31, and 31 $-$ 26 = 5, which is F.
\begin{hints}
\item ``Code A with H, then code the result with K. Where did A go?''
\item ``How many places does H move a letter? K? Both together?''
\item ``T then M moves 31 places. Where is that on a ring of 26?''
\end{hints}

\pb{7}{\core}{Write the rule for the number of pieces and explain it.}
\ans Pieces = the biggest number that goes into both the number of dots and the hop.
\textbf{An explanation a child can give on the 12-dot ring with hop 8.} Number the dots 0 to 11. The first piece
lands on 0, 8, 4 and back to 0. Those are all in the 4-times table, because 8 and 12 are, and going past 12 just
takes 12 away. So dots 1, 2, 3 are never reached and each needs its own start: at least 4 pieces. Each piece is
the first one turned one dot along, so every piece has the same number of dots. The pencil is back when the hops
add up to whole laps: 8, 16, 24, and 24 is two laps after 3 hops. So a piece has 3 dots, and 12 dots make
12 $\div$ 3 = 4 pieces. The same two steps work for any ring and hop.
\begin{hints}
\item ``Make a table of dots, hop and pieces from Problems 1--3. What do you notice?''
\item ``Which dots does the first piece land on? Which times table are they in?''
\item ``Why can the first piece never reach dot 1?''
\end{hints}

\pb{8}{\further}{Rings from 4 to 20 dots where every hop draws a line to every dot with one start.}
\ans \textbf{5, 7, 11, 13, 17, 19.}
\textbf{Why the list is complete.} If the number of dots is in a smaller number's times table (other than 1), hop
by that number: you land only on its times table and miss dots. Hop 2 rules out 4, 6, 8, 10, 12, 14, 16, 18, 20,
and hop 3 rules out 9 and 15. For 5, 7, 11, 13, 17, 19: the pieces all have the same number of dots, so they share
the dots out in equal piles. Seven counters make equal piles only as one pile of 7 or seven piles of 1, and a
piece of one dot means the hop never moved. So every hop gives one piece.
\begin{hints}
\item ``Try 6: which hop misses dots? Now try 7.''
\item ``For a ring that fails, how is the failing hop related to the number of dots?''
\item ``Can 7 counters go into equal piles, with more than one pile and more than one in each?''
\end{hints}

\pb{9}{\further}{Erased dots: find the number of dots and a hop.}
\ans Top left \textbf{15 dots, hop 6 or 9} (three five-pointed stars). Top right \textbf{20 dots, hop 5 or 15}
(five squares). Bottom left \textbf{14 dots, hop 6 or 8} (two seven-pointed stars). Bottom right \textbf{21 dots,
hop 6 or 15} (three seven-pointed stars). Either hop is right: hop $k$ and hop $n-k$ draw the same lines.
\begin{hints}
\item ``Every dot was the end of a line. Count the points around the outside.''
\item ``Follow one line from a point to its other end. Counting around the outside, how many points along is that?''
\item ``Count the separate shapes and the corners of one. Does shapes $\times$ corners give the dots?''
\end{hints}

\pb{10}{\further}{Is coding twice harder to crack?}
\ans No. Two settings in a row are one setting (H then K is R, Problem 6), so Lee's code is an ordinary
one-setting code and still falls in at most 25 tries. If his two settings add to a full turn, such as H then T,
the message comes out not coded at all.
\begin{hints}
\item ``Look at Problem 6. What did H then K make?''
\item ``How many different single settings are there?''
\end{hints}

\pb{11}{\further}{Setting C from A; which settings take A through all 26 letters?}
\ans Setting C: A, C, E, G, I, K, M, O, Q, S, U, W, Y, then A again (13 letters). All 26:
\textbf{B, D, F, H, J, L, P, R, T, V, X, Z} (12 settings).
\textbf{Why the list is complete.} The letters are a ring of 26 dots and a setting is a hop. An even setting
(C, E, \dots, Y) lands only on A, C, E, \dots{} and never on B. N goes A, N, A. A does not move. Every other
setting is odd and not N; the only numbers that go into 26 are 1, 2, 13, 26, so the biggest number going into
both is 1: one piece, all 26 letters.
\begin{hints}
\item ``List the letters for C. Which never appear?''
\item ``Try B, D and E. What do the ones that fail have in common?''
\item ``The wheel is a ring of 26 dots. Use your rule from Problem 7.''
\end{hints}

\pb{12}{\ext}{Times tables on 24 numbered dots: $m \to 2m$ and $m \to 3m$, counting round past 23.}
\ans $\times2$: a heart shape (cardioid) with its dent at dot 0, at the top. Dot 0 gets no line; 8--16 is drawn
twice (8 $\to$ 16 and 16 $\to$ 8), so there are 22 different lines. $\times3$: a two-lobed shape (nephroid) with
dents at dots 0 and 12. Dots 0 and 12 get no line; 3--9, 6--18 and 15--21 are each drawn twice, so there are
19 different lines.
\begin{hints}
\item ``Where does dot 12 go? Dot 13?''
\item ``Draw from dot 1, 2, 3, \dots{} in order and watch the curve appear.''
\end{hints}

% =====================================================================================
\section{The mathematics behind the week}

\textbf{Model.} Number the dots $0,\dots,n-1$ clockwise. A hop of $k$ is $i \mapsto i+k \pmod n$. A start at $s$
draws $s, s+k, s+2k,\dots$ until it returns to $s$.

\textbf{Pieces.} Let $g=\gcd(n,k)$. Then every piece has $n/g$ dots and there are $g$ pieces.
\emph{Proof.} $s+jk\equiv s \pmod n$ iff $n \mid jk$ iff $\tfrac{n}{g} \mid j\tfrac{k}{g}$ iff $\tfrac{n}{g}\mid j$,
since $\gcd(\tfrac{n}{g},\tfrac{k}{g})=1$. So each piece has $n/g$ dots and the $n$ dots split into $g$ pieces.
The piece through 0 is the set of multiples of $k$ mod $n$, which is the set of multiples of $g$ (Bézout:
$g=ak+bn$); the others are its turns by $1,\dots,g-1$ dots. Each piece is the regular star polygon
$\{m/j\}$ with $m=n/g$, $j=k/g$: a polygon when $j\equiv\pm1 \pmod m$, a segment when $m=2$.

\textbf{Hop $k$ and hop $n-k$} draw the same lines in opposite directions, since $-k\equiv n-k$.

\textbf{One piece for every hop iff $n$ is prime.} If $n=ab$ with $1<a<n$, hop $a$ gives $a$ pieces. If $n$ is
prime, $\gcd(n,k)=1$ for $1\le k<n$. In general $\varphi(n)$ hops give one piece, making $\varphi(n)/2$ different
pictures (12 dots: hops 1, 5, 7, 11; two pictures).

\textbf{The wheel is $\mathbb{Z}/26$.} With A $=0,\dots,$ Z $=25$, setting $s$ codes $x\mapsto x+s \pmod{26}$:
hop $s$ on a ring of 26. Two settings make one, $(x+s)+t=x+(s+t)$, and $26-s$ undoes $s$. The 26 shifts form a
group, so coding twice never leaves them (Problem 10), and any code falls in at most 25 tries. Shifts keep word
lengths and repeated letters, which is why BALLOON and one-letter words give a code away.

\textbf{Full cycles.} Repeating setting $s$ from A visits $26/\gcd(26,s)$ letters, by the pieces fact with
$n=26$. All 26 iff $\gcd(s,26)=1$, that is $s$ odd and $s\neq13$: $\varphi(26)=12$ settings, the generators of
$\mathbb{Z}/26$. Setting C ($s=2$) gives 13 letters.

\textbf{Times tables.} As $n$ grows, the chords $m\mapsto km \pmod n$ outline an epicycloid with $k-1$ cusps
(the cardioid for $k=2$, the nephroid for $k=3$). The cusps sit at the fixed points, $(k-1)m\equiv0 \pmod n$:
dot 0 for $\times2$ and dots 0 and 12 for $\times3$ on 24 dots.

\textbf{Where it goes next.} Generators of cyclic groups and Euler's $\varphi$; the multiplying code
$x\mapsto ax$ and the affine code $x\mapsto ax+b$, which can be undone exactly when $\gcd(a,26)=1$ (Problem 11
again, seen through multiplication); codes that change setting letter by letter (Vigenère), which defeat the
try-all-settings attack; the order of a perfect shuffle of a deck of cards; Coxeter's star polygons $\{n/k\}$.

% =====================================================================================
\section{After the session: what to write down}

For each table, write which problems the children reached, what they tried and said (their words), and a
drawing or claim worth keeping (photograph the page). For each claim, tick how far it got: \textbf{tried} (one
case), \textbf{guessed} (said without checking), \textbf{checked} in some cases, or \textbf{explained} (a reason
given).

\newcommand{\rec}[1]{%
\par\needspace{14\baselineskip}\vspace{6pt}\noindent\textbf{#1}\par
\noindent Problems reached (circle):}
\newcommand{\claimrow}[1]{#1 & $\square$ & $\square$ & $\square$ & $\square$ \\}
\newcommand{\claimhead}{\begin{tabularx}{\textwidth}{@{}L c c c c@{}}
\toprule Claim to listen for & tried & guessed & checked & explained \\ \midrule}
\newcommand{\notelines}[1]{\par\noindent #1\par\vspace{2pt}\noindent\rule{\textwidth}{0.4pt}\par\vspace{14pt}\noindent\rule{\textwidth}{0.4pt}\par}

\rec{K--1} 1 \quad 2 \quad 3 \quad 4 \quad 5 \quad 6 \quad 7 \quad 8 \quad 9 \quad 10
\par\smallskip
\claimhead
\claimrow{Hop 1 lands on every dot.}
\claimrow{On 6 dots, hop 2 and hop 3 come back too soon.}
\claimrow{On 7 dots every hop lands on every dot.}
\claimrow{A hop and the hop that changes a picture back make 6.}
\claimrow{Other: \hrulefill}
\bottomrule
\end{tabularx}
\notelines{What they tried and said; a drawing worth keeping:}

\rec{Grades 2--3} 1 \quad 2 \quad 3 \quad 4 \quad 5 \quad 6 \quad 7 \quad 8 \quad 9
\par\smallskip
\claimhead
\claimrow{The number of starts is the hop (true only when the hop goes evenly into the dots).}
\claimrow{The number of starts goes evenly into the number of dots.}
\claimrow{Hop 8 on 12 dots is hop 4 the other way.}
\claimrow{16 dots cannot have 3 starts, because 16 does not split into 3 equal piles.}
\claimrow{A setting and its undo add up to 26.}
\bottomrule
\end{tabularx}
\notelines{What they tried and said; a drawing worth keeping:}

\rec{Grades 4--5} 1 \quad 2 \quad 3 \quad 4 \quad 5 \quad 6 \quad 7 \quad 8 \quad 9 \quad 10 \quad 11 \quad 12
\par\smallskip
\claimhead
\claimrow{Pieces = the biggest number that goes into both the dots and the hop.}
\claimrow{A ring has one piece for every hop exactly when its number has no smaller factors (prime).}
\claimrow{Two settings make one: add them, going round past Z.}
\claimrow{The odd settings other than N take A through all 26 letters.}
\claimrow{Coding twice is no harder to crack.}
\bottomrule
\end{tabularx}
\notelines{What they tried and said; a drawing or claim worth keeping:}

\end{document}
